\documentclass[10pt,a4paper]{amsart}
\usepackage[margin=19mm]{geometry}
\usepackage{amsmath,amssymb,mathtools}
\usepackage[scr=rsfs,cal=euler]{mathalfa}
\usepackage{xfrac}
\usepackage[unicode,hidelinks,bookmarks=false]{hyperref}
\newcommand{\ie}{i.e.\ }
\newcommand{\cf}{cf.\ }
\newcommand{\resp}{respectively\ }
\newcommand{\Subsection}{\subsection}
\providecommand{\nobreakdash}{\nobreak\hbox{-}}
\setlength{\emergencystretch}{2em}
\multlinegap0pt
\usepackage{amssymb}
\usepackage{float}
\newtheorem{lemma}{Lemma}
\newtheorem{proposition}{Proposition}
\title{Nonlinear Stability of Taylor-Couette Flows with Heat Buoyancy}
\author{Yeping Li, Gaofeng Wang, Tianfang Wu}
\date{Source: arXiv:2603.22936v1 (CC BY 4.0)}
\begin{document}
\maketitle

\noindent\emph{Source and licence.} Nonlinear Stability of Taylor-Couette Flows with Heat Buoyancy. Version arXiv:2603.22936v1 (CC BY 4.0), distributed by arXiv under the Creative Commons Attribution 4.0 International licence, http://creativecommons.org/licenses/by/4.0/, as recorded in the arXiv licence metadata for this exact version and shown on the abstract page https://arxiv.org/abs/2603.22936v1. Full licence notice: \texttt{licenses/arxiv-2603.22936-CC-BY-4.0.txt}.\par\smallskip

\noindent\textit{Editorial scope.} Counted unit: the article's proposition pro5.2 (source0.tex lines 695--707). The article's complete original proof of it is delivered with it (source0.tex lines 708--777, one \texttt{proof} environment of 70 lines). No mathematics is written, repaired or re-derived: the statement and the proof are the article's own lines, in the article's own notation, and no retained line is substituted or re-flowed.  Retained uncounted as the source-local context the proof needs: lines 440--447; lines 628--636.  Nothing was omitted from any retained span, so the package carries no omission record.  No retained line contains a citation, so the wrapper carries no bibliography and no bibliography entry.  No retained line contains a figure environment, an image-inclusion command or drawing code, so no image is delivered and the package declares no asset dependency.  Editorial formatting only, in addition: an AMS \texttt{amsart} preamble that restates the article's own declarations and macros used by the retained lines, namely the environments \texttt{lemma}, \texttt{proposition} on separate counters, together with the standard packages those lines need.  The retained environments print in this document as Lemma 1, Proposition 1 in order of appearance, whereas the article prints the same items with the numbering of its own full document.  The complete native source of the article as distributed in the arXiv e-print for this exact version is bundled unchanged as \texttt{sources/197045/source0.tex}.
\begin{lemma}[Weighted elliptic estimates for azimuthal modes] \label{lem2.1}
For $|k| \geq 1$, let $\varphi$ satisfy $\left(\partial_r^2 - \frac{k^2 - 1/4}{r^2}\right) \varphi = \omega$ with $\varphi|_{r=1,R} = 0$. The following inequalities hold:
\begin{align}
\|\partial_r \varphi\|_{L^2}^2 + |k|^2 \|r^{-1}\varphi\|_{L^2}^2 &\lesssim |\langle \omega, \varphi\rangle| \lesssim |k|^{-2} \|r \omega\|_{L^2}^2,\\
\|\partial_r \varphi\|_{L^2}^2 + |k|^2 \|r^{-1}\varphi\|_{L^2}^2 &\lesssim |k|^{-1} \|r^{1/2} \omega\|_{L^1}^2,\\
\|r^{1/2} \partial_r \varphi\|_{L^{\infty}} + |k| \|r^{-1/2} \varphi\|_{L^{\infty}} &\lesssim \left(\frac{R}{R-1}\right)^{1/2} |k|^{-1/2} \|r \omega\|_{L^2}.
\end{align}
\end{lemma}

\begin{equation}\label{5.1}
    \left\{
    \begin{aligned}
    &\partial_t \omega_k+\mathcal{L}_\nu \omega_k=h_1-g\partial_{r}h_2, \\
    &\left(\partial_r^2-\frac{k^2-\frac{1}{4}}{r^2}\right)  \varphi_k=\omega_k,\\
    &\left. (\omega_k,\varphi_k)\right|_{r=1, R}=0, \quad \left. \omega_k\right|_{t=0}=\omega_k(0).\\
\end{aligned}
    \right.
\end{equation}

\begin{proposition} \label{pro5.2}
    For $k \in \mathbb{Z} \backslash\{0\}$ and  $\nu k^2\geq |B|$. Let $\omega_k$ be the solution to \eqref{5.1} with $\omega_k(0) \in L^2$, then there exists a constant $c^{\prime}>0$ independent of $\nu, B, k, R$, such that it holds
    $$
    \begin{aligned}
        &   \left\|\mathcal{E}_k(\omega_k)\right\|_{L^\infty L^2}
        +\nu^{\frac{1}{2}}\left\|\mathcal{E}_k(\partial_{r}\omega_k)\right\|_{L^2L^2}\\
  &+(\nu k^2)^{\frac{1}{2}}\left\|\mathcal{E}_k(\frac{\omega_k}{r})\right\|_{L^2L^2}\\
  & \quad+(\nu k^2)^{\frac{1}{2}} R^{-2}\left(|k|\left\|\mathcal{E}_k(\partial_r \varphi_k)\right\|_{L^2 L^2}+k^2\left\|\mathcal{E}_k(\frac{\varphi_k}{r})\right\|_{L^2 L^2}\right) \\
        &\leq C\left\|\omega_k(0)\right\|_{L^2}+\nu^{-\frac{1}{2}}\left\|\mathcal{E}_k(\frac{rh_1}{k})\right\|_{L^2 L^2}\\
        &\quad+\nu^{-\frac{1}{2}}\left\|\mathcal{E}_k(|g|+r|\partial_{r}g|)h_2\right\|_{L^2L^2}.
    \end{aligned}
    $$
\end{proposition}

\begin{proof}
By performing integration by parts on the above expression and
taking its real part, one gets
$$
\begin{aligned}
& \text{Re}\left\langle\partial_t \omega_k-\nu\left(\partial_r^2-\frac{k^2-\frac{1}{4}}{r^2}\right) \omega_k+\frac{i k B}{r^2} \omega_k-h_1+g\partial_{r}h_2, w_k\right\rangle \\
&\quad =\frac{1}{2} \partial_t\left\|\omega_k\right\|_{L^2}^2+\nu\left\|\partial_r \omega_k\right\|_{L^2}^2+\nu\left(k^2-\frac{1}{4}\right)\left\|\frac{\omega_k}{r}\right\|_{L^2}^2\\
&\quad-\text{Re}\left\langle \frac{rh_1}{k}, k \frac{\omega_k}{r}\right\rangle-\text{Re}\left\langle h_2,  \partial_r\left(g \omega_k\right)\right\rangle=0 .
\end{aligned}
$$
It infers
\begin{equation}
\begin{aligned}
   & \partial_t\|\omega_k\|_{L^2}^2+\nu\left\|\partial_r \omega_k\right\|_{L^2}^2+\nu k^2\left\|\frac{\omega_k}{r}\right\|_{L^2}^2 \lesssim\left\|\frac{rh_1}{k}\right\|_{L^2}|k|\left\|\frac{\omega_k}{r}\right\|_{L^2}\\
    &+\left\|rh_2\partial_rg\right\|_{L^2}\left\|\frac{\omega_k}{r}\right\|_{L^2}+\left\|h_2g\right\|_{L^2}\left\|\partial_r\omega_k\right\|_{L^2} .
    \end{aligned}
\end{equation}
By applying Cauchy-Schwarz inequality, we then deduce
$$
\partial_t\left\|\omega_k\right\|_{L^2}^2+\nu\left\|\partial_r \omega_k\right\|_{L^2}^2+\nu k^2\left\|\frac{\omega_k}{r}\right\|_{L^2}^2 \lesssim \nu^{-1}\left(\left\|\frac{rh_1}{k}\right\|_{L^2}^2+\left\|rh_2\partial_rg\right\|_{L^2}^2+\left\|h_2g\right\|_{L^2}^2\right) .
$$
Noticing that $\left\|\dfrac{\omega_k}{r}\right\|_{L^2} \geq R^{-1}\left\|\omega_k\right\|_{L^2}$, it follows
\begin{equation}
\begin{aligned}
&\partial_t\left\|\omega_k\right\|_{L^2}^2+\nu\left\|\partial_r \omega_k\right\|_{L^2}^2+\nu k^2 R^{-2}\left\|\omega_k\right\|_{L^2}^2+\nu k^2\left\|\frac{\omega_k}{r}\right\|_{L^2}^2 \\
&\lesssim \nu^{-1}\left(\left\|\frac{rh_1}{k}\right\|_{L^2}^2+\left\|rh_2\partial_rg\right\|_{L^2}^2+\left\|h_2g\right\|_{L^2}^2\right) .
\end{aligned}
\end{equation}
Therefore, due to $ \nu k^2R^{-2} \geq( \nu
k^2)^{1/3}|B|^{2/3}R^{-2}$,  we can multiply $e^{c^{\prime}\left(\nu
k^2\right)^{\frac{1}{3}}|B|^{\frac{2}{3}} R^{-2} t}$ on both sides
of above inequality. With $c^{\prime}$ being a small constant
independent of $\nu, B, k, R$, we can obtain
$$
\begin{aligned}
& \partial_t\left\|\mathcal{E}_k (\omega_k)\right\|_{L^2}^2+\nu\left\|\mathcal{E}_k (\partial_r \omega_k)\right\|_{L^2}^2+\nu k^2\left\|\mathcal{E}_k (\frac{\omega_k}{r})\right\|_{L^2}^2 \\
&\lesssim \nu^{-1}\left\|\mathcal{E}_k (\frac{rh_1}{k})\right\|_{L^2}^2+\nu^{-1}\left\|\mathcal{E}_k(rh_2\partial_rg)\right\|_{L^2}^2\\
&\quad+\nu^{-1}\left\|\mathcal{E}_k(
h_2g)\right\|_{L^2}^2,
\end{aligned}
$$
which further implies
$$
\begin{aligned}
& \left\|\mathcal{E}_k (\omega_k)\right\|_{L^{\infty} L^2}^2
+\nu\left\|\mathcal{E}_k (\partial_r \omega_k)\right\|_{L^2 L^2}^2\\
&+\nu k^2\left\|\mathcal{E}_k (\frac{\omega_k}{r})\right\|_{L^2 L^2}^2 \\
& \lesssim\left\|\omega_k(0)\right\|_{L^2}^2
+\nu^{-1}\left\|\mathcal{E}_k(\frac{rh_1}{k})\right\|_{L^2 L^2}^2
+\nu^{-1}\left\|\mathcal{E}_k (rh_2\partial_rg)\right\|_{L^2 L^2}^2\\
&\quad+\nu^{-1}\left\|\mathcal{E}_k (h_2g)\right\|_{L^2L^2}^2 .
\end{aligned}
$$
In view of Lemma \ref{lem2.1}, it also holds that
$$
R^{-2}\left(|k|\left\|\varphi_k^{\prime}\right\|_{L^2}+k^2\left\|\frac{\varphi_k}{r}\right\|_{L^2}\right) \lesssim R^{-2}\left\|r \omega_k\right\|_{L^2} \leq\left\|\frac{\omega_k}{r}\right\|_{L^2}.
$$
Combining these two estimates above yields
$$
\begin{aligned}
& \left\|\mathcal{E}_k (\omega_k)\right\|_{L^{\infty} L^2}+\nu^{\frac{1}{2}}\left\|\mathcal{E}_k (\partial_r \omega_k)\right\|_{L^2 L^2}\\
&\quad+\left(\nu k^2\right)^{\frac{1}{2}}\left\|\mathcal{E}_k (\frac{\omega_k}{r})\right\|_{L^2 L^2} \\
&+\left(\nu k^2\right)^{\frac{1}{2}} R^{-2}\left(|k|\left\|\mathcal{E}_k (\varphi_k^{\prime})\right\|_{L^2}+k^2\left\|\mathcal{E}_k(\frac{\varphi_k}{r})\right\|_{L^2}\right) \\
& \lesssim\left\|\omega_k(0)\right\|_{L^2}
+\nu^{-\frac{1}{2}}\left\|\mathcal{E}_k (\frac{rh_1}{k})\right\|_{L^2 L^2}+\nu^{-\frac{1}{2}}\left\|\mathcal{E}_k (rh_2\partial_rg)\right\|_{L^2 L^2}\\
&\quad+\nu^{-\frac{1}{2}}\left\|\mathcal{E}_k (h_2g)\right\|_{L^2L^2},
\end{aligned}
$$
which completes the proof.
\end{proof}

\end{document}
